\chapter{Introdução}
\label{ch:introducao}

A organização multi-escala é traço distintivo dos sistemas vivos. As células
coordenam transformações bioquímicas em escala de milissegundos por meio de
redes metabólicas, ao mesmo tempo que controlam programas de desenvolvimento
em escala de horas ou dias por meio de redes regulatórias. Decisões celulares
como o crescimento versus a diferenciação, a quiescência versus a proliferação e
a adaptação versus a decisão irreversível emergem da interação entre
essas duas escalas operacionais. A predição quantitativa dessas decisões ---
em especial dos \emph{limiares de decisão}, fronteiras quantitativas nas quais
as células transitam de forma irreversível entre fenótipos --- permanece desafio
central da biologia de sistemas computacional.

A motivação imediata para este trabalho provém de medições experimentais precisas
da esporulação de \textit{Bacillus subtilis}. \textcite{Fujita2005}
demonstraram que a acumulação gradual de Spo0A${\sim}$P via fosforretransmissão
--- KinA~$\to$~Spo0F~$\to$~Spo0B~$\to$~Spo0A~$\to$~Spo0A${\sim}$P ---
produz esporulação em ${\approx}\,52\%$ das células, enquanto a indução
constitutiva abrupta de Spo0A* produz apenas ${\approx}\,5\%$. A separatriz
de comprometimento emerge quando $\sigma^H$ ultrapassa
$\theta_{\sigma H} \approx 1{,}60\,\mu$M, concentração que delimita o
cruzamento irreversível para o programa de diferenciação. Essa assimetria
suscita questões fundamentais: por que a via gradual supera tão amplamente a
indução abrupta? Essa fronteira de comprometimento é
determinada pela topologia da rede regulatória, pela cinética das enzimas, pelos
limiares das vias de sinalização ou por alguma combinação desses fatores? Qual
o formalismo adequado para integrar informações metabolômicas e regulatórias de
modo a calcular tais limites a partir da topologia da rede e de parâmetros cinéticos mensuráveis?

\section{Motivação: o Problema das Escalas Cruzadas}
\label{sec:motivacao}

As Redes de Petri~\cite{Petri1962,Murata1989} estabeleceram-se como formalismo
canônico para modelagem de sistemas com dinâmicas concorrentes e discretas.
Extensões posteriores --- redes coloridas, temporizadas, estocásticas e
híbridas --- alargaram o poder expressivo para acomodar tokens tipados, semânticas
temporais, disparo probabilístico e marcações de valor real. Contudo, uma lacuna
persistente separa modelagem metabólica (fluxo de massa estequiométrico, cinética
de Michaelis-Menten, redes de ocorrências) de modelagem regulatória (lógica
Booleana, equações diferenciais ordinárias, lógica temporal). Os sistemas
biológicos requerem simultaneamente as duas semânticas, e as abordagens
existentes obrigam a uma escolha arquitetural que muitas vezes obscurece
precisamente os fenômenos de interesse.

Essa limitação manifesta-se concretamente nas moléculas com papéis duplos.
O ATP funciona como substrato da glicólise (semântica de transferência de massa,
contabilidade estequiométrica) e como sinal de controle para cascatas de
fosforilação (semântica de transferência de informação, detecção de limiar e
consumo de sinal). O cálcio participa da energética da contração muscular e
desencadeia cascatas de transdução de sinal. O AMP cíclico acopla o metabolismo
de carbono à repressão por catabólito. Forçar essas moléculas a uma única
semântica introduz inconsistências formais que impossibilitam predições
quantitativas. A biologia de sistemas necessita de um formalismo capaz de
expressar a dualidade metabolito-sinal de forma matematicamente coerente.

\section{Nota Metodológica}
\label{sec:nota-metodologica}

O objetivo deste trabalho é \textbf{validação de formalismo}, não descoberta
biológica. A biologia da esporulação de \textit{B.~subtilis} é bem estabelecida;
os parâmetros do modelo provêm de dados publicados; e o resultado esperado ---
a assimetria entre acumulação gradual de Spo0A${\sim}$P
(${\approx}\,52\%$ de esporulação) e indução constitutiva abrupta de Spo0A*
(${\approx}\,5\%$)~\cite{Fujita2005} --- é conhecido antecipadamente.
Esse enquadramento é intencional: um formalismo confiável deve reproduzir
fenômenos conhecidos antes de ser utilizado para predizer fenômenos novos. O
critério de êxito não é, portanto, ``descobrimos algo novo sobre
\textit{B.~subtilis}?'', mas sim ``o formalismo prediz o fenômeno a partir da
topologia da rede e de parâmetros cinéticos, sem ajuste ao resultado?''. A reprodução pela simulação
da separatriz emergente de $\sigma^H$ a $\theta_{\sigma H} \approx 1{,}60\,\mu$M
--- que discrimina os destinos celulares gradual/abrupto sem ajuste de
nenhum parâmetro ao dado de Fujita --- satisfaz esse critério.

\section{A Teoria das Redes de Petri Hierárquicas de Sinais}
\label{sec:teoria-shpn}

Para abordar essas questões, este trabalho apresenta as \textbf{Redes de Petri
Hierárquicas de Sinais} (\textit{Signal Hierarchical Petri Nets}, SHPN), uma
extensão formal das Redes de Petri clássicas que acolhe transferência de massa
e fluxo de informação regulatória dentro de uma estrutura matemática unificada.
O formalismo distingue \emph{lugares de sinal} --- espécies moleculares que
intermediam controle regulatório além de participar no metabolismo --- e
introduz \emph{arcos de fluxo de sinal} que consomem tokens durante a
avaliação de habilitação, criando fronteiras de decisão irreversíveis.
Uma \emph{função de camada hierárquica} organiza os sinais por escala
(metabolismo, sinalização, transcrição), garantindo que a depleção de camadas
inferiores preempta as superiores~\cite{Simao2025HierarchicalPreemption}.
O parâmetro termodinâmico $\Gamma = (K, n, \varepsilon)$, composto por
constantes cinéticas mensuráveis, fundamenta os limiares de habilitação sem
recorrer a ajuste experimental.
A definição formal completa e as propriedades teóricas são apresentadas no
Capítulo~\ref{ch:shpn}.

\section{Limitações das Abordagens Clássicas}
\label{sec:limitacoes-classicas}

As abordagens existentes de biologia de sistemas computacional abordam a
modelagem metabólico-regulatória com diferentes graus de expressividade, mas
nenhuma oferece o conjunto completo de garantias necessárias à predição de
limiares de decisão.

As Redes de Petri clássicas e coloridas modelam com excelência a estequiometria
metabólica e a concorrência discreta, mas carecem de semântica hierárquica para
expressar controle regulatório de nível superior sobre processos metabólicos.
Os modelos ODE (equações diferenciais ordinárias) capturam dinâmicas contínuas
com fidelidade cinética, porém não representam explicitamente a topologia da
rede: decisões qualitativas emergem das equações, mas nunca são verificáveis como
propriedades formais. O SBML (\textit{Systems Biology Markup Language})~\cite{Hucka2003}
define formato de intercâmbio, mas não formalismo de análise; a semântica
regulatória depende do simulador subjacente. As redes Booleanas capturam a lógica
regulatória de forma compacta, mas eliminam a cinética enzimática e a
estequiometria metabólica, impossibilitando predições quantitativas precisas.
Os processos estocásticos~\cite{Gillespie1977,Goss1998} tratam o ruído
molecular, mas a variabilidade estocástica nos limiares de decisão não se deriva
formalmente da topologia da rede. Por fim, as redes híbridas~\cite{David1992}
combinam lugares contínuos e discretos mas exigem dois paradigmas de modelagem
distintos, reintroduzindo a separação metabólico-regulatória que se pretende
suprimir.

A lacuna analítica é precisa: nenhum desses formalismos separa formalmente
transferência de massa de transferência de informação, nenhum aplica semântica
de consumo a arcos regulatórios, e nenhum deriva limiares de decisão a
partir da estrutura da rede sem ajuste de parâmetros.

\section{Roteiro do Trabalho}
\label{sec:roteiro}

A tese está organizada em quatro partes. A Parte~I (Fundamentos, Capítulos~1
a~3) apresenta as bases teóricas: este capítulo enquadra o problema, o
Capítulo~2 revisa os formalismos existentes e identifica as lacunas, e o
Capítulo~3 introduz a definição formal completa das SHPN com semântica
operacional, o parâmetro termodinâmico $\Gamma$ e resultados de complexidade.
A Parte~II (Validação, Capítulo~4) obtém validação quantitativa com o
modelo de esporulação de \textit{B.~subtilis}, demonstrando que o limiar de
decisão emerge da dinâmica termodinâmica de $\Gamma$. A Parte~III
(Implementação, Capítulo~5) descreve a plataforma SHYpn --- implementação
computacional de referência do formalismo SHPN --- e suas integrações com as
bases KEGG e BRENDA.
Por fim, a Parte~IV (Síntese, Capítulos~6 e~7) contextualiza o
formalismo no panorama da teoria de Redes de Petri, discute questões em aberto
e enuncia as contribuições.

Leitores com interesse primordialmente teórico podem percorrer os Capítulos~1,
2 e~3 em sequência e depois consultar a síntese do Capítulo~6. Aqueles
interessados em aplicações biológicas podem prosseguir da introdução diretamente
para o Capítulo~4, retomando a fundamentação formal conforme necessário.
Pesquisadores de biologia de sistemas computacional encontrarão no Capítulo~5 a
descrição do fluxo de trabalho de modelagem automatizada integrada às bases de
dados.
